AlphaFold2 prédit la structure tridimensionnelle d'une protéine à partir de sa seule séquence, en exploitant les signaux de coévolution contenus dans un alignement multiple de séquences homologues. Ce projet met en place un workflow reproductible \emph{séquence $\rightarrow$ AlphaFold2 $\rightarrow$ structure $\rightarrow$ analyse}. Après comparaison de plusieurs solutions, nous avons retenu ColabFold~1.5.5 exécuté localement dans Docker : il utilise les mêmes poids qu'AlphaFold2, mais remplace les 2,6~To de bases de données par un serveur MMseqs2, et fonctionne sans GPU. Le workflow a été testé sur la calmoduline humaine (UniProt P0DP23, 149 résidus), prédite en 1~h~57 sur CPU. Le pLDDT (moyenne 82,8) montre deux lobes bien prédits reliés par une région centrale moins fiable. La PAE révèle que, si chaque lobe est rigide (environ 4--5~Å), leur orientation relative reste incertaine (environ 20~Å). Un \texttt{README.md} et deux scripts permettent de refaire une prédiction et son analyse en deux commandes.
